\chapter[Metodologia]{Metodologia}

Este Capítulo estabelece a estrutura e as técnicas empregadas para desenvolver e validar o modelo de previsão, garantindo a replicabilidade e a solidez científica do projeto. A metodologia está organizada de forma sequencial e detalha as seguintes etapas: a aquisição e o pré-processamento dos dados, a modelagem, o treinamento, os objetivos e métricas de desempenho, as arquiteturas de \textit{Deep Learning}, as formas de avaliação e a utilização de métodos de Interpretabilidade xAI.


\section{Escopo do Projeto}

Foi montado com o objetivo de transformar dados históricos de preço e indicadores de mercado em previsões acionáveis para o período futuro. A estratégia adotada utiliza a LSTM como modelo \textit{baseline} para conduzir uma avaliação sistemática de diferentes abordagens de pré-processamento, engenharia de \textit{features} e representações de \textit{input}. Após esse estudo, são realizadas previsões no mesmo período de tempo (2000--2024) em arquiteturas de \textit{Deep Learning} SOTA (xLSTM-TS e Transformer-TS), visando verificar se modelos mais complexos conseguem extrapolar o desempenho do \textit{baseline} e obter previsões direcionais mais precisas para o índice S\&P~500. A Figura \ref{fig:fluxo_processamento} ilustra o fluxo de trabalho e orienta a organização das seções seguintes.

\begin{itemize}
    \item \textbf{Aquisição de Dados e Pré-processamento:} Coleta de preços históricos e volume do S\&P~500 (2000--2024) e tratamento, com definição dos períodos de treino, validação e teste.
    \item \textbf{\textit{Ablation Study} (LSTM \textit{Baseline}):} Construção de \textit{features} e teste de seis variantes de \textit{input} e normalização; avaliação via MCC, Acurácia, F1, AUC e \textit{McNemar}.
    \item \textbf{Modelos SOTA:} Realização das previsões com o mesmo período do \textit{ablation study} nos modelos xLSTM-TS e Transformer-TS; avaliação via métricas de regressão (RMSE, MAPE, $R^2$).
    \item \textbf{Direção Derivada dos Modelos SOTA:} Obtém uma previsão por classificação comparando o valor da previsão com o valor do índice no dia, indicando se a tendência é de aumentar ou diminuir o valor.
    \item \textbf{Comparação de Desempenho dos Modelos:} Avaliação do desempenho dos modelos de acordo com as métricas de classificação (Acurácia e F1).
    \item \textbf{Interpretabilidade xAI:} Aplicação do SHAP para identificar a importância temporal de cada \textit{timestep} da janela de entrada.
\end{itemize}

\begin{figure}[H]
    \centering
    \includegraphics[height=0.86\textheight]{figuras/flow.png}
    \caption{Fluxo de projeto para previsão direcional do S\&P500, da coleta de dados ao pipeline de xAI.}
    \label{fig:fluxo_processamento}
\end{figure}
% \clearpage

% Diagrama criado no link:<https://www.mermaidchart.com/app/projects/5aed441e-b401-4aee-b551-99767edd6e57/diagrams/28745ebb-4d79-4843-bcb9-d9f85363a7a4/version/v0.1/edit>
\section{Aquisição de Dados}

Esta seção apresenta as informações referentes ao conjunto de dados utilizado no estudo, incluindo as fontes de informação, o período de observação considerado e as variáveis empregadas como entradas para a modelagem preditiva. O foco deste trabalho é o mercado acionário norte-americano, escolhido em razão de sua elevada liquidez e relativa estabilidade quando comparado ao mercado brasileiro. Para a validação do modelo proposto, será utilizado como referência o Índice S\&P~500, indicador de desempenho do mercado de ações dos Estados Unidos.

Os dados históricos serão coletados via Yahoo Finance \cite{site17} com granularidade diária, cobrindo o período de 03 de janeiro de 2000 a 30 de dezembro de 2024, totalizando 6.288 dias úteis.  O \textit{dataset} inclui as variáveis padrão de negociação (OHLC) e volume, base das análises técnicas e cálculo das \textit{features}.

\section{Pré-processamento de Dados}

A fase de pré-processamento é essencial para transformar a série temporal despadronizada em um formato estruturado, capaz de maximizar a capacidade de aprendizado dos modelos de previsão. Esta etapa prioriza a escalabilidade dos dados e a criação de uma estrutura temporal que evite o vazamento de informações futuras. As escolhas a seguir foram realizadas em estudos como o de \citeonline{Gupta2025Interpretable}, que demonstram a eficácia dessa abordagem em contextos semelhantes de previsão de séries temporais financeiras.

\begin{itemize}
    \item \textbf{Tratamento de Ruído:}

    Antes de construir qualquer \textit{feature}, os dados brutos foram inspecionados para garantir sua integridade. A série do S\&P~500 (2000--2024) não apresentou valores faltantes nas colunas de preço (OHLC). O único problema identificado foram dias com volume zero, que foram tratados como ausentes e imputados pela mediana móvel dos dias anteriores. Além disso, foi definido que o preço de referência utilizado em todo o projeto seria o \textit{Close} (preço de fechamento bruto), em vez do \textit{Adj Close} (preço de fechamento ajustado), por ser o valor bruto de mercado sem ajustes retroativos que poderiam introduzir viés temporal.

    \item \textbf{Engenharia de \textit{Features}:}

    Os preços brutos (OHLCV) não são diretamente adequados como entrada para redes neurais em séries temporais financeiras. O preço absoluto do S\&P~500 cresceu de $\approx$1.400~USD em 2000 para $\approx$6.000~USD em 2024, o que significa que os valores do conjunto de teste estão muito fora da escala do treino. Para contornar isso, as \textit{features} foram construídas como quantidades relativas e adimensionais derivadas do preço e do volume, que descrevem o comportamento do mercado em vez do seu nível absoluto.

    As 18 \textit{features} utilizadas podem ser agrupadas em quatro categorias e estão detalhadas a seguir:

    \begin{itemize}
        \item \textbf{Log-retornos e defasagens} (7 \textit{features}): o log-retorno diário atual ($\log(P_t/P_{t-1})$) mede a variação percentual do preço de forma simétrica e estacionária (Preço $P_t$ é o \textit{Close} do dia $t$). As defasagens de 1, 2, 3, 5, 10 e 20 dias fornecem ao modelo memória sobre os movimentos recentes do mercado, totalizando 7 \textit{features} de retorno.
        \item \textbf{Tendência relativa} (3 \textit{features}): razões entre o preço atual e as médias móveis de 5 e 20 dias ($P_t / \overline{P}_5$, $P_t / \overline{P}_{20}$, $\overline{P}_5 / \overline{P}_{20}$). A média móvel $\overline{P}_k$ é calculada como a média dos preços de fechamento dos últimos $k$ dias. Essas razões indicam se o mercado está acima ou abaixo da sua tendência de curto e médio prazo.
        \item \textbf{Volatilidade realizada} (2 \textit{features}): desvio padrão dos log-retornos em janelas de 20 e 60 dias, calculado como $\sigma_k = \sqrt{\frac{1}{k}\sum_{i=0}^{k-1} (\textit{logret}_{t-i} - \bar{\textit{logret}})^2}$, onde $\textit{logret}_{t-i} = \log(P_{t-i}/P_{t-i-1})$. Capturam o regime de risco do mercado: período calmo x período de alta incerteza.
        \item \textbf{Volume e amplitude} (6 \textit{features}): as seis \textit{features} deste grupo são calculadas da seguinte forma: (i) volume em escala logarítmica, definido como $\log(V_t)$ onde $V_t$ é o volume negociado no dia $t$; (ii) volume normalizado pela média de 20 dias, calculado como $V_t / \overline{V}_{20}$, onde $\overline{V}_{20}$ é a média dos volumes dos últimos 20 dias; (iii) amplitude intradiária, dada por $H_t - L_t$, sendo $H_t$ o preço máximo e $L_t$ o preço mínimo do dia $t$; (iv) retorno absoluto, definido como $|P_t - P_{t-1}|$; (v) retorno quadrático, calculado como $(P_t - P_{t-1})^2$; e (vi) RSI de 14 períodos, índice de força relativa calculado como $\textit{RSI}_t = 100 - \frac{100}{1 + \textit{RS}_t}$, onde $\textit{RS}_t = \frac{\overline{G}_{14}}{\overline{L}_{14}}$ é a razão entre a média dos ganhos ($\overline{G}_{14}$) e a média das perdas ($\overline{L}_{14}$) dos últimos 14 dias. O RSI varia de 0 a 100, com valores acima de 70 indicando sobrecompra e abaixo de 30 indicando sobrevenda. Essas \textit{features} refletem a intensidade das negociações e o momentum do mercado.
    \end{itemize}

    O \textit{target} é a variável que o modelo precisa prever, funcionando como a resposta correta que o algoritmo deve aprender a identificar. Para cada dia $t$, o target $y_{t+1}$ indica a direção do preço no dia seguinte: 1 se o preço subir, 0 se cair. O target é sempre calculado a partir do preço \textit{Close} original, sem aplicar nenhum pré-processamento (como DWT ou normalização). Isso evita vazamento de informação: o rótulo de saída permanece puro e independente das transformações aplicadas às \textit{features} de entrada.

    \item \textbf{\textit{Ablation Study} --- Variantes de Normalização e \textit{Input}:}

    Para identificar a representação de \textit{input} e a normalização mais adequadas para os modelos de previsão, foram definidas seis variantes que cobrem sistematicamente três eixos de decisão: (i) tipo de \textit{input} do preço; (ii) aplicação ou não de \textit{denoising wavelet}; e (iii) técnica de normalização. Todos os parâmetros de normalização e limiares de \textit{denoising} são ajustados exclusivamente nos dados de treino, garantindo que nenhuma informação do conjunto de validação ou teste influencie o pré-processamento. A Tab.~\ref{tab:variantes_norm} apresenta as seis variantes testadas e as motivações de cada uma.

    \begin{table}[H]
    \centering
    \caption{Variantes do \textit{Ablation Study} de Pré-processamento}
    \label{tab:variantes_norm}
    \resizebox{\textwidth}{!}{%
    \begin{tabular}{|c|c|c|p{0.55\textwidth}|}
    \hline
    \textbf{Variante} & \textbf{\textit{Input}} & \textbf{Normalização} & \textbf{Motivação} \\
    \hline
    A & \textit{Close} bruto & Z-score & \textit{Baseline} padrão da literatura. As 18 \textit{features} são calculadas a partir do \textit{Close} original, normalizado por média e desvio padrão do treino. Ponto de referência mínimo. \\
    \hline
    B & \textit{Close} DWT & Z-score & Aplica DWT no \textit{Close} antes de calcular as 18 \textit{features} e o Z-score. Avalia se remover o ruído de alta frequência do preço melhora a aprendizagem do modelo. \\
    \hline
    C & Log-retornos & Z-score & Substitui o \textit{Close} pelo log-retorno diário como base para calcular as 18 \textit{features}. Avalia a previsão do modelo em uma série estacionária e adimensional. \\
    \hline
    D & Log-ret. DWT & Z-score & Aplica DWT nos log-returns antes de calcular as 18 \textit{features}. Combina a estacionarização via log-retorno com o DWT. Avalia se ambas as transformações juntas são complementares. \\
    \hline
    E & \textit{Close} bruto & Min-Max & As 18 \textit{features} são calculadas a partir do \textit{Close} original, normalizado para o intervalo $[0, 1]$. Testa a sensibilidade do modelo a uma normalização que depende dos valores extremos do treino. \\
    \hline
    F & \textit{Close} DWT & Min-Max & Aplica DWT no \textit{Close} antes de calcular as 18 \textit{features} e a normalização Min-Max. Combina o DWT com a normalização Min-Max. \\
    \hline
    \end{tabular}%
    }
    \end{table}

    \item \textbf{Divisão Temporal e Treinamento:}

    A divisão foi feita de forma estritamente temporal, sem embaralhamento, para simular condições reais de implantação. O modelo utiliza uma janela deslizante de 20 dias (\textit{lookback}) para prever a direção do preço no dia seguinte (\textit{timestep} $t+1$):
    \begin{itemize}
        \item \textbf{Treino:} 2000--2020 ($\approx$5.275 dias, 84\%) --- inclui crises de 2008 e COVID-2020, permitindo ao modelo aprender múltiplos regimes econômicos.
        \item \textbf{Validação:} 2021--2022 ($\approx$503 dias, 8\%) --- monitoramento do \textit{early stopping} e ajuste de hiperparâmetros.
        \item \textbf{Teste:} 2023--2024 ($\approx$503 dias, 8\%) --- regime de juros altos pós-pandemia, contexto macroeconômico inédito para o modelo.
    \end{itemize}

    Para as utilização do \textit{denoising} via Transformada Wavelet Discreta (DWT), o limiar de corte foi calculado exclusivamente nos dados de treino. Essa normalização também foi aplicada nos períodos de validação e teste apenas para avaliação das previsões, evitando vazamento de informações do futuro para o passado, pois o modelo de treino nunca tem acesso às estatísticas dos períodos subsequentes.

    A validação cruzada \textit{K-Fold} foi descartada por ser inadequada para séries temporais: na divisão aleatória, dados futuros inevitavelmente aparecem no treino, gerando métricas artificialmente otimistas. O método \textit{Walk Forward} também não foi adotado, pois na prática exige um treinamento antes de cada previsão, incluindo o dia anterior previsto, o que aumentaria significativamente o tempo de processamento devido à necessidade de múltiplos treinamentos e previsões. Como o objetivo do projeto é comparar diferentes arquiteturas de \textit{Deep Learning}, optou-se pelo treinamento simples com \textit{Time Series Split} (divisão temporal fixa), que simula condições reais de implantação e evita vazamento de dados.

\end{itemize}

\section{Modelagem e Arquitetura}

Foram implementados três modelos de \textit{Deep Learning} com diferentes abordagens de previsão. A LSTM foi utilizada como \textit{baseline} para o \textit{ablation study}; para os modelos SOTA (xLSTM-TS e Transformer-TS) foram realizados alguns ajustes ao pipeline original para garantir causalidade e evitar vazamento de dados: (i)~substituição do modo de DWT de \texttt{per} (periódico) para \texttt{zero} (causal), evitando vazamento de informação futura; (ii)~ajuste do MinMaxScaler apenas no treino (sem vazamento); e (iii)~derivação da classificação direcional a partir da diferença entre previsões consecutivas de regressão.

\begin{itemize}    
    \item \textbf{LSTM (Classificação Binária Direta):}

    A LSTM foi implementada com Keras/TensorFlow como classificador binário e utilizada como \textit{baseline} para o \textit{ablation study} por possuir tempo de processamento bem inferior às demais arquiteturas. O modelo consiste em uma camada LSTM unidirecional com 32 unidades, \textit{dropout} de 0{,}1 e uma camada densa de saída com ativação sigmoide que produz $P(\textit{sobe}) \in [0,1]$. A função de perda é a Entropia Cruzada Binária (\textit{Binary Crossentropy}) e o otimizador é o Adam com taxa de aprendizado padrão, \textit{batch size} de 32 e até 100 épocas. O \textit{input} são as 18 \textit{features} técnicas sobre uma janela de 20 dias (\textit{lookback}). O treinamento utiliza \textit{early stopping} com paciência de 10 épocas, monitorando a perda de validação e restaurando os melhores pesos.


    \item \textbf{xLSTM-TS (Regressão com Direção Derivada):}

    O xLSTM-TS \citeonline{Gil2024Evaluation} foi implementado em PyTorch a partir do repositório oficial, com as seguintes adaptações para este trabalho: (i)~substituição do \textit{backend} CUDA pelo \textit{backend vanilla} para execução em CPU; (ii)~redução do \textit{input} para uma única \textit{feature} (\textit{Close} denoised), em vez das múltiplas séries utilizadas no repositório original; (iii)~escalonamento Min-Max ajustado exclusivamente nos dados de treino, evitando vazamento de informação. O pré-processamento aplicado foi: \textit{denoising} DWT causal (\textit{db4}, \texttt{mode='zero'}, limiar calculado no treino) seguido de normalização Min-Max, com janela de entrada (\textit{lookback}) de 150 dias. Os demais hiperparâmetros do artigo original foram preservados: 4 blocos xLSTM, \textit{embedding} de dimensão~64, 2 cabeças de atenção, \textit{batch size}~16, taxa de aprendizado $10^{-4}$, agendamento \textit{ReduceLROnPlateau} e \textit{early stopping} com paciência de 40 épocas. O modelo é treinado para regressão do preço absoluto (perda MSE) e a classificação de subida ou descida é derivada em pós-processamento pela diferença entre o valor absoluto da previsão e o valor do índice no dia.


    \item \textbf{Transformer-TS (Regressão com Direção Derivada):}

    O Transformer-TS foi implementado em PyTorch a partir da arquitetura proposta por \citeonline{Vaswani2017Attention}, com adaptações para regressão de séries temporais univariadas: utilizou-se apenas o bloco codificador (sem o decodificador original) e adicionou-se uma cabeça de regressão sobre o último passo temporal. A arquitetura resultante consiste em projeção de entrada Linear$(1 \to 64)$, \textit{positional embedding} aprendido, 2 blocos \textit{TransformerEncoderLayer} com atenção de 4 cabeças e FFN de dimensão~256, totalizando $\approx$113.921 parâmetros treináveis. O pré-processamento e todos os hiperparâmetros de treino foram alinhados aos do xLSTM-TS (mesma perda MSE, otimizador, \textit{scheduler}, \textit{early stopping} e \textit{lookback} de 150 dias), garantindo comparação direta entre as duas arquiteturas.
\end{itemize}

Neste trabalho, o Transformer-TS e o xLSTM-TS são comparados sob condições idênticas de treinamento, permitindo avaliar o ganho marginal de cada arquitetura SOTA em relação à LSTM \textit{baseline}.

\subsection{Ferramentas e Implementação}

Para garantir a reprodutibilidade dos experimentos, todas as implementações utilizaram as seguintes versões de frameworks e bibliotecas: Python 3.10, TensorFlow 2.15 (para LSTM baseline), PyTorch 2.3.1 (para xLSTM-TS e Transformer-TS) e SHAP 0.46. Ambos os modelos SOTA foram treinados para regressão do preço absoluto com perda MSE e com hiperparâmetros idênticos (4 blocos, embedding 64, 2 cabeças, batch size 16, lr $10^{-4}$, \textit{early stopping} com paciência 40). Os modelos SOTA foram baseados em implementações de código aberto: xLSTM-TS foi adaptado do repositório oficial de \citeonline{Gil2024Evaluation} e o Transformer-TS foi implementado seguindo a arquitetura descrita em \citeonline{Vaswani2017Attention}. O código-fonte está disponível em repositório público para fins acadêmicos em \citeonline{Hansen2026}.

\section{Objetivo das Previsões}

O objetivo principal dos três modelos é a classificação direcional: para cada dia $t$ do conjunto de teste, prever se o preço de fechamento do dia $t+1$ será maior (classe~1 --- sobe) ou menor (classe~0 --- desce). A LSTM realiza essa tarefa diretamente como classificador; o xLSTM-TS e o Transformer-TS a realizam indiretamente, por meio de regressão de preço e derivação da direção pela diferença consecutiva.

As métricas utilizadas foram selecionadas para cobrir diferentes aspectos da avaliação e desempenho:

\begin{itemize}
    \item \textbf{MCC}: adotado como métrica principal de classificação direcional, por considerar desequilíbrios na proporção de dias de alta e baixa.
    \item \textbf{Acurácia, F1-\textit{Score} e AUC-ROC}: métricas complementares de classificação que permitem comparar os modelos sob diferentes perspectivas e facilitar o diálogo com a literatura.
    \item \textbf{RMSE, MAPE e {R}²}: aplicados ao xLSTM-TS e Transformer-TS para avaliar a qualidade das previsões dos modelos de regressão.
    \item \textbf{Teste de \textit{McNemar}}: para verificar se as diferenças de desempenho entre variantes de pré-processamento e arquiteturas são estatisticamente significativas.
\end{itemize}

\section{Estratégia de Interpretabilidade (xAI)}

A última etapa da metodologia consiste em aplicar o SHAP aos modelos treinados, garantindo transparência e auditabilidade das previsões. A análise foi executada com \textit{background} exclusivamente do conjunto de treino e avaliada sobre 200 amostras do conjunto de teste, evitando vazamento de dados. A escolha de 200 amostras equilibra precisão estatística e tempo de processamento, seguindo práticas comuns em estudos de xAI para séries temporais, já que o SHAP possui um custo computacional alto e geralmente é aplicado em um grupo representativo de dados. Os valores SHAP foram calculados via \textit{KernelExplainer} para a LSTM (modelo Keras) e via \textit{DeepExplainer} para os modelos PyTorch (xLSTM-TS e Transformer-TS). 

A lógica do SHAP baseia-se na teoria dos valores de Shapley, que atribui a cada \textit{feature} (neste caso, cada \textit{timestep} da janela de entrada) uma contribuição marginal para a previsão final. Para cada amostra de teste, o SHAP calcula quanto cada \textit{timestep} contribuiu para afastar a previsão do valor base (média das previsões do conjunto de treino). A importância global de cada \textit{timestep} é obtida pela média da magnitude dessas contribuições ($|SHAP|$) sobre todas as amostras analisadas. O foco da análise é:

\begin{itemize}
    \item \textbf{Contribuição Temporal dos \textit{Time Steps}}: o SHAP mede a contribuição de cada dia do histórico de entrada para a previsão final. Isso permite identificar quais períodos da ``memória'' da rede foram mais determinantes para a decisão do modelo.
    \item \textbf{Concentração da Atenção}: a curva de importância acumulada (Pareto) indica quantos \textit{timesteps} são necessários para explicar 90\% da importância total, revelando se o modelo distribui sua atenção ao longo de toda a janela ou a concentra em poucos dias recentes.
\end{itemize}
